%! Factoring Trinomials — Unit 1, Lesson 1.4 — Student Packet Cover Sheet
\documentclass[10pt]{article}
\usepackage{atda-article}
\usepackage{atda-boxes}
\usepackage{ltablex}
\keepXColumns

\begin{document}

% Full-bleed royal banner
\begin{tikzpicture}[remember picture, overlay]
  \fill[royal] (current page.north west)
    rectangle ([yshift=-0.9in]current page.north east);
\end{tikzpicture}
\vspace{-0.8in}

\begin{center}
  {\color{white}\textbf{\LARGE Algebra, Trigonometry, and Data Analysis}}\\[4pt]
  {\color{mistmid}\large\textbf{Unit 1: Algebraic Expressions and Factoring}}\\[6pt]
  {\color{white}\textbf{\large Lesson 1.4 \quad Factoring Trinomials}}
\end{center}

\vspace{0.22in}
\namedateperiod
\vspace{0.18in}

\begin{learningtargetbox}
\small
\begin{enumerate}[leftmargin=1.5em, itemsep=5pt]
  \item \ldots{} factor a trinomial $x^2+bx+c$ by finding two integers whose \textbf{product} is
        $c$ and whose \textbf{sum} is $b$.
  \item \ldots{} factor a trinomial $ax^2+bx+c$ with $a \ne 1$ by \textbf{splitting the middle
        term} --- the $ac$ method --- and then grouping.
  \item \ldots{} factor \textbf{completely}, GCF first, and decide when a trinomial is
        \textbf{prime} over the integers.
  \item \ldots{} read a factored area model to recover the \textbf{dimensions} of the figure it
        describes, and say what those dimensions mean.
\end{enumerate}
\end{learningtargetbox}

\vspace{0.14in}

\begin{tocbox}
\small
\renewcommand{\arraystretch}{1.5}
\begin{tabularx}{\linewidth}{@{} c l X r @{}}
\rowcolor{mist}
\textbf{\#} & \textbf{Component} & \textbf{Description} & \textbf{Score} \\
\midrule
1 & Warm-Up        & Two products, one grouping, and where the split comes from
                                                                    & \blank{1.2cm} \\
2 & Guided Notes   & Product-and-sum, the $ac$ method, and factoring completely
                                                                    & \blank{1.2cm} \\
3 & Group Activity & The cafeteria mural: recovering a shape from its area
                                                                    & \blank{1.2cm} \\
4 & Exit Ticket    & Two trinomials and one answer to critique      & \blank{1.2cm} \\
5 & Homework       & Mixed practice, a garden bed, and a look ahead to special forms
                                                                    & \blank{1.2cm} \\
\midrule
  & \multicolumn{2}{r}{\textbf{Total}} & \blank{1.2cm} \\
\end{tabularx}
\end{tocbox}

\vspace{0.14in}

\begin{remindbox}
\small
This lesson continues \texttt{A2.EO.3b} --- factoring polynomials completely, in one and two
variables, over the integers. Nothing here replaces Lesson 1.3: a trinomial has three terms, and
\textbf{splitting the middle term} turns it into the four-term grouping problem you can already do.
Two habits still matter most: \textbf{GCF first, every time}, and \textbf{check by multiplying
back}.
\end{remindbox}

\end{document}
